\chapter{Conclusões e Trabalhos Futuros}
\label{chap:conclusoes-e-trabalhos-futuros}

O desenvolvimento de veículos aéreos não tripulados, especialmente no contexto de competições de aerodesign, demanda o monitoramento rigoroso de variáveis cinemáticas e atmosféricas em tempo real para a validação de modelos aerodinâmicos e estruturais. Diante do alto custo e da natureza proprietária das soluções comerciais de instrumentação em voo, este trabalho propôs a concepção e implementação de um sistema de telemetria distribuído, estruturado de ponta a ponta, visando prover uma alternativa de baixo custo, alta modularidade e arquitetura aberta.

A síntese das etapas de desenvolvimento demonstra que o objetivo primário foi plenamente alcançado. A arquitetura foi fragmentada em três componentes funcionais integrados: o nó embarcado na aeronave, o middleware terrestre e a estação de controle em solo (Odin). No domínio embarcado, a integração de um microcontrolador de duplo núcleo (ESP32) com sensores inerciais (IMU), barométricos e de posicionamento global (GNSS) garantiu a aquisição quantizada dos dados sob estritas restrições elétricas e de peso. A transmissão dessas leituras operou sobre o protocolo LoRa operante em 915 MHz, estabelecendo um enlace de radiofrequência caracterizado por longo alcance e alta resiliência contra desvanecimentos.

Para o ambiente de solo, a introdução de um middleware dedicado atuou como um conversor eficiente de protocolos físicos, roteando o tráfego do espectro LoRa para uma rede local sem fio em formato de datagramas UDP. Esta decisão arquitetural isolou a estação de operação computacional das complexidades de baixo nível do rádio. Por fim, o software cliente cumpriu os requisitos de monitoramento ao adotar o padrão arquitetural Model-View-Controller (MVC), provendo uma interface orientada a eventos capaz de deserializar, renderizar gráficos de alta frequência e registrar os logs de voo de forma assimétrica, sem incorrer em bloqueios de processamento (thread-blocking).

A consolidação destes componentes, atestada pelos ensaios elétricos e pelas rotinas de simulação Hardware-in-the-Loop (HIL), confirmou a viabilidade técnica do sistema. Com uma taxa de fragmentação de pacotes contida em níveis operacionais aceitáveis (próxima a 1\%) e estabilidade funcional mantida sob condições de estresse na malha de potência, conclui-se que a solução entrega uma plataforma telemétrica robusta, pronta para amparar ensaios de voo práticos com a precisão e a confiabilidade exigidas pelo escopo do projeto.

\section{Contribuições Técnicas e Metodológicas}

Sob a ótica da engenharia de \textit{software} embarcado, a principal contribuição metodológica deste trabalho reside na aplicação da Arquitetura Hexagonal  na concepção do \textit{firmware} da aeronave. Em contraste com abordagens monolíticas comuns em projetos de instrumentação estudantil, a estruturação baseada na injeção de dependências promoveu um alto nível de desacoplamento, isolando rigorosamente a lógica de domínio — responsável pela quantização, escalonamento matemático e construção das máquinas de estados — das amarras físicas de \textit{hardware}. Como resultado direto, essa abstração viabilizou a execução de metodologias de teste em ambiente controlado, como o \textit{Hardware-in-the-Loop}. Esta capacidade de substituir as interfaces sensoriais por adaptadores emulados permitiu a injeção de perfis de voo sintéticos, garantindo a validação exaustiva do sistema e a mitigação de falhas de \textit{software} de forma antecipada e segura.

Do ponto de vista da arquitetura de sistemas e de redes de comunicação, o projeto demonstra mérito ao inserir uma camada de \textit{middleware} terrestre operando como ponte translatora entre a radiofrequência sub-GHz e a rede local de alta capacidade. A delegação da recepção física dos pacotes LoRa e do subsequente encapsulamento em datagramas UDP a um nó microcontrolado isentou o computador da estação de controle (Odin) da severa sobrecarga de processamento associada à escuta contínua de barramentos seriais (UART/USB) em baixo nível.

Esta segmentação arquitetural proporcionou um fluxo de dados limpo e padronizado via \textit{sockets} de rede, permitindo que a aplicação cliente operasse de forma totalmente assíncrona. Consequentemente, o uso desta topologia descentralizada revelou-se uma solução eficaz para evitar gargalos comuns em sistemas de recepção direta, assegurando que o processamento do computador de solo estivesse integralmente dedicado à renderização gráfica a 20 Hz e à gravação redundante dos pacotes MAVLink em disco, sem comprometer a latência ou o acúmulo de dados.


\section{Limitações Encontradas}

A despeito dos resultados satisfatórios alcançados e da validação funcional da arquitetura, o desenvolvimento e a experimentação deste projeto esbarraram em limitações inerentes às restrições de escopo, de instrumentação laboratorial e das próprias tecnologias adotadas. O reconhecimento destas restrições é fundamental para a avaliação crítica do sistema e para o balizamento das expectativas em cenários de operação reais.

A primeira limitação evidente manifestou-se na camada física de comunicação. Embora a taxa de invalidação de pacotes de 1\% ateste a resiliência do protocolo LoRa sob condições adversas, ela confirma a suscetibilidade do enlace sub-GHz a obstruções na primeira Zona de Fresnel e à perda da linha de visada direta (LoS). Em cenários de voo que envolvam a interposição de relevo acidentado denso ou estruturas urbanas entre a aeronave e a base, a degradação do sinal pode escalar, indicando que a solução atual, baseada em antenas omnidirecionais simples, possui um limiar de saturação em ambientes de forte obstrução.

No escopo da validação elétrica do sistema embarcado, identificou-se uma limitação técnica significativa relacionada aos equipamentos de medição empregados. Consequentemente, o aparato experimental foi incapaz de capturar quedas de tensão transientes de altíssima frequência que podem ser induzidas pelo acionamento abrupto dos servomotores da aeronave. Embora o circuito supervisor interno do microcontrolador não tenha registrado reinicializações, a ausência de análises com osciloscópios de alta velocidade impede a garantia de imunidade a transientes eletromagnéticos severos.

Por fim, no que tange à infraestrutura física, a montagem estrutural do componente \textit{middleware} apresenta restrições mecânicas. A decisão de instanciar o circuito de recepção terrestre sobre \textit{protoboards} com condutores rígidos foi justificada pela imobilidade da estação e pela facilidade de manutenção em bancada. Contudo, essa topologia carece da robustez mecânica e da blindagem contra interferências eletromagnéticas (EMI) inerentes a uma Placa de Circuito Impresso (PCB) dedicada. Em um ambiente de pista de voo, a exposição a vibrações acidentais ou manuseio contínuo eleva o risco de mau contato nos terminais, o que poderia comprometer temporariamente o roteamento dos pacotes para a rede local da estação de controle.


\section{Trabalhos Futuros}

Como consequência natural do ciclo de desenvolvimento e visando o amadurecimento contínuo da arquitetura proposta, elenca-se um conjunto de diretrizes para a evolução do sistema. 

No escopo do \textit{hardware}, recomenda-se a transição dos protótipos experimentais para Placas de Circuito Impresso dedicadas a nível de \textit{Surface Mount Device} (SMD), abrangendo tanto o nó embarcado na aeronave quanto o componente de \textit{middleware} terrestre. Esta evolução construtiva é necessária para garantir a miniaturização física, a redução da massa estrutural e, criticamente, a mitigação de interferências eletromagnéticas em malhas de alta frequência. Ademais, a base sensorial da aeronave pode ser expandida através da integração de instrumentos específicos de dinâmica de fluidos, como um tubo de Pitot acoplado a um transdutor de pressão diferencial, viabilizando a medição direta da velocidade aerodinâmica relativa.

No que tange ao \textit{software}, sugere-se uma exploração mais aprofundada das capacidades operacionais da camada de \textit{middleware}. O microcontrolador de solo pode ter sua responsabilidade expandida para atuar além de um roteador passivo de protocolos. A implementação de lógicas de processamento de borda (\textit{edge computing}) no \textit{middleware} permitiria a pré-filtragem de pacotes corrompidos, a análise da qualidade de serviço (QoS) do enlace LoRa em tempo real e a agregação de estatísticas de conexão antes da retransmissão dos datagramas via UDP, otimizando o consumo de banda na rede local da estação.

Para o aprimoramento dos métodos de verificação, propõe-se a adoção de processos formais de qualidade de \textit{software} e de metodologias consolidadas de engenharia. A implementação de rotinas de testes unitários focadas estritamente na camada de domínio garantirá a corretude matemática contínua das rotinas de quantização e conversão de escalas. Em paralelo, a adoção da Engenharia de Sistemas Baseada em Modelos (MBSE) proverá uma estrutura analítica rigorosa para a rastreabilidade de requisitos, gestão de interfaces e controle do ciclo de vida do projeto, mitigando o risco de regressões arquiteturais em iterações futuras do código.

Por fim, a infraestrutura de comunicação concebida pode ser expandida para unificar todo o ecossistema de dados em voo da equipe. A arquitetura modular desenvolvida permite que outras frentes de instrumentação, como estações portáteis de aquisição de dados climáticos para aferição de densidade do ar na pista, sejam integradas de forma transparente à plataforma unificada. Ao padronizar a saída desses novos nós sensores para o protocolo MAVLink e roteá-los através da malha UDP existente, a estação de solo será capaz de correlacionar variáveis ambientais exógenas com o desempenho cinemático da aeronave, enriquecendo profundamente a capacidade de análise pós-voo sem demandar reestruturações profundas no \textit{software} base.
